\documentclass[a4paper,10pt]{article}
\usepackage[margin=18mm]{geometry}
\usepackage{fontspec}
\setmainfont[Scale=0.92]{Noto Sans}
\usepackage{polyglossia}\setmainlanguage{russian}
\usepackage{tabularx,booktabs,array}
\newcommand{\field}[1]{\textbf{#1}}
\pagestyle{empty}
\begin{document}
\begin{center}\Large\bfseries Есеп-төлем парағы / Расчётный листок\\
\normalsize {{ period_label }}\end{center}
\vspace{4mm}
\noindent\field{Жұмыс беруші / Работодатель:} {{ employer.name }}, БСН/БИН {{ employer.bin }}\\
\noindent\field{Қызметкер / Работник:} {{ employee.full_name }}, ЖСН/ИИН {{ employee.iin }}\\
\noindent\field{Лауазымы / Должность:} {{ employee.position }}
\vspace{4mm}
\noindent\begin{tabularx}{\textwidth}{@{}X r@{}}
\toprule
Көрсеткіш / Показатель & Сомасы, ₸ / Сумма, ₸ \\
\midrule
Есептелген жалақы (жалпы сома) / Начислено (гросс) & {{ gross_tenge }} \\
\midrule
Жеке табыс салығы (ИПН) / Индивидуальный подоходный налог (ИПН) & {{ ipn }} \\
Міндетті зейнетақы жарналары (ОПВ) / Обязательные пенсионные взносы (ОПВ) & {{ opv }} \\
Жұмыскердің ӘМСС жарнасы (ВОСМС) / Взнос работника на ОСМС (ВОСМС) & {{ vosms }} \\
\midrule
\bfseries Қолға берілетін сома / К выплате (нетто) & \bfseries {{ net }} \\
\midrule
Жұмыс беруші есебінен зейнетақы жарналары (ОПВР) / Взносы работодателя в ЕНПФ (ОПВР) & {{ opvr }} \\
Әлеуметтік аударым (СО) / Социальные отчисления (СО) & {{ so }} \\
Жұмыс берушінің ӘМСС жарнасы (ОСМС) / Взнос работодателя на ОСМС (ОСМС) & {{ osms }} \\
Әлеуметтік салық / Социальный налог & {{ social_tax }} \\
\bottomrule
\end{tabularx}
\vspace{3mm}
\noindent Қолға берілетін сома: {{ net }} ₸\\
({{ net_words }})
\vspace{12mm}
\noindent Бас бухгалтер / Гл. бухгалтер \hrulefill\qquad Қызметкер / Работник \hrulefill
\end{document}
